\documentclass[12pt,a4paper]{article}

% --- Packages ---
\usepackage{amsmath,amssymb}
\usepackage{graphicx}
\usepackage{booktabs}
\usepackage{natbib}
\usepackage[margin=1in]{geometry}
\usepackage{setspace}
\usepackage{hyperref}
\usepackage{xcolor}
\usepackage{float}
\usepackage{threeparttable}
\usepackage{tabularx}
\usepackage{array}
\usepackage{enumitem}
\usepackage{caption}

\hypersetup{colorlinks=true, linkcolor=blue, citecolor=blue, urlcolor=blue}
\doublespacing

% --- Title ---
\title{The True Cost of Volatility Targeting: Decomposing the Insurance Premium into Opportunity and Transaction Components\thanks{All errors are my own. Data and replication code are available upon request.}}

\author{Yi-Hao Lai\thanks{Department of Finance, Da-Yeh University. Corresponding author. Email: yhlai@mail.dyu.edu.tw}}

\date{April 2026}

\begin{document}

\maketitle

% --- Abstract ---
\begin{abstract}
\noindent
Volatility targeting (VT) strategies scale equity exposure inversely to expected volatility, producing well-documented reductions in tail risk but also persistent return shortfalls relative to buy-and-hold. We propose a decomposition of this return gap---the ``insurance premium''---into two economically distinct components: \emph{opportunity cost} (foregone returns from reduced equity participation) and \emph{direct cost} (transaction expenses from weight rebalancing). Using daily S\&P~500 data conditioned on the CBOE VVIX index over the VVIX-reliable period 2012--2024 ($N=3{,}262$ trading days), we find that opportunity cost accounts for 91\% of the total insurance premium under unconditional VT (4.20\% vs.\ 0.43\% per annum). VVIX-conditional targeting---activating VT only when the volatility-of-volatility $z$-score exceeds 1.0---reduces total cost by 74\% in the full sample (from 4.62\% to 1.22\%), primarily by eliminating unnecessary hedging during calm markets. However, cross-OOS analysis over six non-overlapping windows shows that this cost reduction is sample-dependent: S2 outperforms buy-and-hold on a Sharpe-ratio basis in only 2 of 6 windows. We therefore treat VVIX conditioning as a hypothesis-generating in-sample accounting result rather than a validated trading rule. We further show that the 50/50 SPY/GLD benchmark generates a full-sample average rebalancing premium of 54~basis points per annum (negative in two of four sub-periods), raising the long-run bar for VT to clear. While VoV-conditional VT (S2) achieves a higher Sharpe ratio than 50/50 (0.63 vs.\ 0.50), the comparison is not apples-to-apples: S2 is 100\% equity with occasional VT, while 50/50 includes gold diversification. These findings redirect the VT design problem from minimizing transaction costs to measuring when opportunity cost, rather than direct trading friction, dominates the insurance premium.

\medskip
\noindent
\textbf{Keywords:} volatility targeting, insurance premium, opportunity cost, transaction cost, volatility-of-volatility, rebalancing premium

\medskip
\noindent
\textbf{JEL Classification:} G11, G12, G17
\end{abstract}

\newpage

%=================================================================
\section{Introduction}
\label{sec:intro}
%=================================================================

Volatility targeting (VT) has become a standard technique for dynamic risk management. The framework, in the influential formulation of \citet{moreira2017}, prescribes scaling equity exposure inversely to conditional volatility---increasing allocation when markets are calm and reducing it when turbulence rises. \citet{harvey2018} demonstrate that such strategies meaningfully reduce tail risk; \citet{fleming2001} establish the economic value of volatility timing, and \citet{hocquard2013} document thinner left tails under constant-volatility designs.

Yet VT strategies systematically underperform buy-and-hold benchmarks in terms of cumulative returns. This shortfall is commonly attributed to trading friction: implementation costs feature prominently in practitioner discussions of VT, and \citet{cederburg2020} document poor out-of-sample performance across 103 volatility-managed strategies. However, this attribution conflates two fundamentally different economic mechanisms. The first is the \emph{opportunity cost} of holding cash instead of equities during periods of elevated volatility that do not materialize into realized losses. The second is the \emph{direct cost} of portfolio turnover---brokerage fees, bid-ask spreads, and market impact. Despite a substantial literature on VT mechanics \citep{moreira2017, barroso2015, cederburg2020, bongaerts2020, liu2019}, the decomposition of this insurance premium into its constituent components has received little empirical attention.

While \citet{bongaerts2020} propose conditional volatility targeting by conditioning on volatility states to optimize Sharpe ratios, our approach conditions on volatility-of-volatility and focuses on cost decomposition rather than performance optimization. This distinction matters for at least three reasons. First, our decomposition clarifies the economic nature of VT: investors are paying primarily for foregone upside, not for trading friction. Second, it implies that cost reduction should target the opportunity component rather than the transaction component, fundamentally reorienting strategy design. Third, it explains a persistent empirical puzzle: why does VT fail to outperform simple 50/50 SPY/GLD diversification, even when VT produces superior risk metrics?

We answer this question with one primary contribution and two supporting empirical exercises. The primary contribution is a simple accounting decomposition of the VT insurance premium into opportunity and direct cost components (Section~\ref{sec:method}). While the decomposition is arithmetically straightforward, its empirical magnitude---91\% opportunity cost at conservative friction assumptions---has not been previously documented and challenges the prevailing focus on transaction cost minimization. As supporting evidence, we quantify the full-sample average rebalancing premium of 54~basis points embedded in the 50/50 benchmark (with material sub-period variation, including two negative sub-periods), highlighting the long-run bar that VT must clear to outperform simple diversification. We also report an exploratory VVIX-conditioning exercise that reduces the total premium by 74\% in the full sample, primarily through the opportunity cost channel, but whose cross-OOS success rate remains weak at 2 of 6 windows.

%=================================================================
\section{Methodology and Data}
\label{sec:method}
%=================================================================

\subsection{Insurance Premium Decomposition}
\label{sec:decomp}

Let $r_t^{\text{BH}}$ denote the daily return of the buy-and-hold equity benchmark and $r_t^{\text{VT}}$ the daily return of the volatility-targeted strategy. The VT strategy allocates weight $w_t = \min(\sigma^*/\hat{\sigma}_t, 1)$ to equities, where $\sigma^*$ is the target volatility and $\hat{\sigma}_t$ is the conditional volatility estimate available at time $t-1$. Following the target-volatility convention studied by \citet{perchet2016}, we implement a 12\% annual volatility target using lagged VIX as the forecast: $w_t = \min(12/\text{VIX}_{t-1}, 1)$.

Define the \emph{total insurance premium} as the annualized return difference, where $\bar{r}$ denotes the annualized mean daily return:
\begin{equation}
\text{IP}_{\text{total}} = \bar{r}^{\text{BH}} - \bar{r}^{\text{VT}}.
\label{eq:ip_total}
\end{equation}

We apply a simple accounting decomposition into two components. Let $r_t^{\text{VT,gross}}$ denote the VT return \emph{before} transaction costs are deducted. The \emph{opportunity cost} captures the return gap due to reduced equity participation:
\begin{equation}
\text{IP}_{\text{opp}} = \bar{r}^{\text{BH}} - \bar{r}^{\text{VT,gross}}.
\label{eq:ip_opp}
\end{equation}

The \emph{direct cost} captures the cumulative transaction expense:
\begin{equation}
\text{IP}_{\text{direct}} = \Bigl( \sum_{t} c \cdot |w_t - w_{t-1}| \Bigr) \Big/ Y,
\label{eq:ip_direct}
\end{equation}
where $c = 5$~basis points per unit weight change and $Y = T/252$ is the sample length in years ($T = 3{,}262$ trading days, $Y \approx 12.94$). By construction, $\text{IP}_{\text{total}} = \text{IP}_{\text{opp}} + \text{IP}_{\text{direct}}$.

\subsection{VoV-Conditional Targeting}

Rather than applying VT unconditionally, we propose conditioning on the volatility-of-volatility (VoV) regime. Using the CBOE VVIX index \citep{cboe2012}, we compute an expanding-window $z$-score (minimum 60 observations):
\begin{equation}
z_t^{\text{VoV}} = \frac{\text{VVIX}_t - \overline{\text{VVIX}}_{1:t}}{\text{sd}(\text{VVIX}_{1:t})}.
\end{equation}

We define two conditional strategies:
\begin{itemize}
\item \textbf{S2 (VoV-Conditional)}: Apply 12/VIX weight only when $z_t^{\text{VoV}} > 1.0$ \emph{and} VIX is rising (5-day change $ > 0$); otherwise hold 100\% equity. The 5-day window balances responsiveness against noise: shorter windows (1--2 days) produce excessive whipsaw, while longer windows (10+ days) delay regime detection. Sensitivity to the threshold $z$-score is reported in Section~\ref{sec:results}.
\item \textbf{S3 (Smooth VoV)}: Interpolate continuously between full equity and full VT as a linear function of the clipped VoV $z$-score: $w_t = 1 - \operatorname{clip}(z_t^{\text{VoV}}, 0, 1)\,(1 - w_t^{\text{VT}})$, equivalently $w_t = w_t^{\text{VT}} + (1 - w_t^{\text{VT}})\bigl(1 - \operatorname{clip}(z_t^{\text{VoV}}, 0, 1)\bigr)$. Thus $z_t^{\text{VoV}} \le 0$ holds 100\% equity, $z_t^{\text{VoV}} \ge 1$ applies full VT ($w_t = w_t^{\text{VT}}$), and intermediate values interpolate linearly. This continuous specification is the one used to produce the S3 results reported below.
\end{itemize}

The economic intuition is that high VoV signals regime uncertainty---periods when insurance is most valuable---while low VoV indicates stable volatility regimes where VT hedging is unnecessary and costly.\footnote{Unreported results show that the VoV $z$-score threshold alone accounts for approximately 80\% of the cost reduction, with the VIX-rising filter contributing the remaining 20\% by preventing premature activation during declining-volatility episodes. We acknowledge that the 80/20 attribution has not been validated through nested out-of-sample testing. It is presented as a directional estimate rather than a precise decomposition.}

\subsection{Data}

We use daily closing prices for SPY, GLD, VIX, and VVIX from Yahoo Finance (via \texttt{yfinance}), consistent with CRSP to within rounding precision, covering the VVIX-reliable period 2012-01-03 to 2024-12-31 ($N = 3{,}262$ trading days, 12.94 years). The VVIX restriction to post-2012 data follows from documented liquidity concerns in the VVIX options market prior to this date. All signals are lagged one trading day (\texttt{signal.shift(1)}) to prevent look-ahead bias. Following a conservative assumption, transaction costs are set at 5~bps per leg on each weight change, which comfortably exceeds the small effective trading costs documented for liquid U.S.\ equities \citep{hasbrouck2009} and is well above SPY's quoted bid-ask spread---on the order of 1--2~bps for this exceptionally liquid ETF---thereby providing an upper bound on the direct cost component. At 1~bps, the opportunity cost share rises to 98.0\%, further strengthening our core finding. We designate the out-of-sample (OOS) window as 2023-01-01 to 2024-12-31.

%=================================================================
\section{Results}
\label{sec:results}
%=================================================================

\subsection{Strategy Performance}

Table~\ref{tab:performance} reports performance metrics for five strategies over the full sample (2012--2024). The buy-and-hold benchmark (S0) achieves a CAGR of 12.51\% with a Sharpe ratio of 0.60. Unconditional VT (S1, Always 12/VIX) sacrifices 5.40 percentage points of CAGR for a 55\% reduction in maximum drawdown (MDD: $-15.46\%$ vs.\ $-34.10\%$). The VoV-conditional strategy (S2) retains most of the BH return (CAGR 11.14\%) while still reducing MDD by 33\% and achieving the highest Sharpe ratio (0.63) among all strategies.

\begin{table}[H]
\centering
\begin{threeparttable}
\caption{Strategy Performance: Full Sample (2012--2024)}
\label{tab:performance}
\begin{tabular}{lccccc}
\toprule
 & S0: BH & S1: Always & S2: VoV- & S3: Smooth & S4: 50/50 \\
Metric & SPY & 12/VIX & Cond. & VoV & SPY/GLD \\
\midrule
CAGR (\%) & 12.51 & 7.11 & 11.14 & 8.84 & 7.89 \\
Ann.\ Vol (\%) & 16.56 & 9.33 & 13.72 & 11.57 & 11.47 \\
Sharpe & 0.60 & 0.54 & 0.63 & 0.57 & 0.50 \\
MDD (\%) & $-$34.10 & $-$15.46 & $-$22.78 & $-$23.20 & $-$21.17 \\
Calmar & 0.37 & 0.46 & 0.49 & 0.38 & 0.37 \\
Sortino & 0.55 & 0.49 & 0.61 & 0.53 & 0.46 \\
Mean CRRA $\gamma\!=\!5$ & 0.217 & 0.190 & 0.211 & 0.208 & 0.157 \\
\midrule
Avg.\ Weight (\%) & 100.0 & 74.1 & 93.4 & 86.0 & --- \\
Ann.\ Turnover (\%) & 0.0 & 872.4 & 1044.7 & 912.4 & --- \\
\bottomrule
\end{tabular}
\begin{tablenotes}
\small
\item \textit{Notes}: S0 is 100\% SPY buy-and-hold. S1 sets $w = \min(12/\text{VIX}_{t-1}, 1)$ daily. S2 applies S1 only when VoV $z$-score $> 1.0$ and VIX is rising; otherwise $w = 1$. S3 continuously interpolates S1 toward full equity as the clipped VoV $z$-score falls below 1.0 (linear in $\operatorname{clip}(z,0,1)$). S4 is a 50/50 SPY/GLD portfolio held at constant daily weights (continuous, costless rebalancing); its Sharpe (0.50) is a daily constant-weight benchmark and is distinct from the separately computed monthly-rebalancing premium (2006--2024) analysed in Section~\ref{sec:results}. All signals lagged one day. TX cost: 5~bps per weight change. CRRA $\gamma\!=\!5$ reports mean CRRA utility of the wealth path (not certainty-equivalent); a formal CE analysis requires computing $\text{CE} = u^{-1}(E[u(W)])$. Data source: yfinance; experiment K811v2.
\end{tablenotes}
\end{threeparttable}
\end{table}

\subsection{Insurance Premium Decomposition}

Table~\ref{tab:decomposition} presents the central result: the decomposition of the insurance premium into opportunity and direct cost components.

\begin{table}[H]
\centering
\begin{threeparttable}
\caption{Insurance Premium Decomposition (\%/year)}
\label{tab:decomposition}
\begin{tabular}{lccccc}
\toprule
Strategy & Opportunity & Direct & Total & Opp.\ Share & $\Delta$ vs.\ S1 \\
 & Cost & Cost & Premium & (\%) & (\%) \\
\midrule
S1: Always 12/VIX & 4.20 & 0.43 & 4.62 & 90.7 & --- \\
S2: VoV-Conditional & 0.70 & 0.52 & 1.22 & 57.1 & $-$73.6 \\
S3: Smooth VoV & 2.85 & 0.46 & 3.31 & 86.2 & $-$28.4 \\
\bottomrule
\end{tabular}
\begin{tablenotes}
\small
\item \textit{Notes}: Opportunity cost $= \bar{r}^{\text{BH}} - \bar{r}^{\text{VT,gross}}$ (return gap before TX). Direct cost $= \sum c|w_t - w_{t-1}|$ annualized (TX friction). Total $=$ Opp.\ $+$ Direct. $\Delta$ vs.\ S1 shows percentage reduction in total premium relative to unconditional VT. Data: K811v2 (2012--2024, $N = 3{,}262$).
\end{tablenotes}
\end{threeparttable}
\end{table}

Three findings emerge. First, for unconditional VT (S1), opportunity cost dominates at 91\% of the total premium (4.20\% vs.\ 0.43\%). This overturns the common practitioner intuition that VT is expensive because of trading costs. The actual cost is foregone equity participation during elevated-VIX periods that are followed by positive returns.

Second, VoV conditioning (S2) reduces the total premium by 74\%, from 4.62\% to 1.22\% per annum. Critically, this reduction operates almost entirely through the opportunity cost channel: S2's opportunity cost falls from 4.20\% to 0.70\% (an 83\% reduction), while its direct cost actually \emph{increases} slightly from 0.43\% to 0.52\%. The mechanism is intuitive: S2 is fully invested 86\% of the time ($w = 1$), engaging VT only during the 14\% of days when VoV $z$-score exceeds 1.0 and VIX is rising. This eliminates the wasteful hedging during calm periods that accounts for the bulk of S1's opportunity cost.

Third, the smooth blending approach (S3) offers a middle ground but is less effective than the binary VoV switch, reducing total cost by only 28\%. This suggests that partial hedging during low-VoV regimes still incurs substantial opportunity cost without proportionate risk reduction.

\subsection{Structural Advantages of the 50/50 Benchmark}

A persistent puzzle is that no VT strategy reliably outperforms a simple 50/50 SPY/GLD portfolio on a Sharpe-ratio basis, despite VT's superior drawdown characteristics. We resolve this puzzle by identifying three structural advantages embedded in the 50/50 benchmark.

First, SPY and GLD exhibit near-zero correlation ($\rho = 0.057$) over the extended 2006--2024 sample period.\footnote{The correlation and rebalancing premium estimates use the 2006--2024 sample (from GLD inception) rather than the 2012--2024 VVIX-reliable period.} Over the 2012--2024 VVIX-reliable period reported in Table~\ref{tab:performance}, this low correlation produces a 50/50 portfolio volatility of 11.47\%---30.7\% below SPY's standalone volatility of 16.56\%. Second, monthly rebalancing of the 50/50 portfolio generates a full-sample average premium of 54~basis points per annum (CAGR of 10.02\% vs.\ 9.49\% for buy-and-hold)---though with substantial sub-period variation: across four sub-periods the empirical premium is negative in two (2006--09: $-95$~bps; 2010--14: $-56$~bps) and positive in two, so it is a long-run average property rather than a per-period guarantee---consistent with the theoretical prediction of \citet{booth1992} and confirmed by our weight-sweep analysis across allocations from 30/70 to 70/30.\footnote{The 54~bps estimate uses dividend-adjusted price series (yfinance \texttt{auto\_adjust=True}, which incorporates dividend reinvestment into the price level), matching the original implementation. The companion replication package (under the \texttt{auto\_adjust=False} raw-Close convention mandated by the project's reproduction rule set) yields approximately 63~bps for the same 2006--2024 window---both estimates fall within the structural $\sim$50--81~bps range reported in the next paragraph and reflect the same buy-low/sell-high mechanism under different dividend-handling conventions.} This premium arises mechanically from systematic buy-low/sell-high behavior when rebalancing to fixed weights. Third, gold provides crisis alpha: during the four VoV regimes identified by VVIX, gold consistently appreciates during equity drawdowns---the ``HighVoV\_Rising'' regime where VT provides its greatest benefit is precisely the regime where gold delivers positive returns, creating a natural hedge that requires no signal, no transaction, and no parameter estimation.

These three structural advantages combine to generate approximately 54--81~bps per annum (the rebalancing premium alone plus the diversification benefit). While S2 achieves a higher Sharpe ratio than 50/50 (0.63 vs.\ 0.50), this reflects fundamentally different risk exposures: S2 is predominantly equity with state-contingent hedging, whereas 50/50 derives its risk-adjusted return from cross-asset diversification and mechanical rebalancing. Comparing them on Sharpe alone obscures this distinction.

\subsection{Regime Analysis}

The insurance premium varies dramatically across VoV regimes. During ``HighVoV\_Falling'' periods (7.7\% of the sample), S1's opportunity cost spikes to 19.95\%---these are the post-crisis recoveries where VT misses the rebound. In contrast, during ``LowVoV\_Falling'' periods (45.5\% of the sample), S1's opportunity cost is a more modest 2.54\%. The VoV-conditional strategy S2 eliminates this regime asymmetry by construction: it holds 100\% equity during low-VoV regimes, incurring zero opportunity cost in the 76\% of days when VoV is below the threshold.

\subsection{Cross-OOS Robustness}

We evaluate robustness using six non-overlapping two-year windows (2013--14, 2015--16, 2017--18, 2019--20, 2021--22, and 2023--24), comparing per-window Sharpe ratios. S2 outperforms BH SPY on a Sharpe-ratio basis in only 2 of 6 windows (2017--18 and 2019--20), while S3 also wins in 2 of 6 and the 50/50 SPY/GLD benchmark wins in 3 of 6. Diebold-Mariano tests\footnote{DM tests use the negative daily return as the loss function ($L_t = -r_t$, so the loss differential is $d_t = r_{2,t} - r_{1,t}$), with Bartlett-kernel Newey--West HAC standard errors and maximum lag $\lceil T^{1/3} \rceil$, matching the archived implementation (\texttt{strategy\_dm\_test}, K811v2).% audit fix 2026-06-10: footnote previously described a QLIKE/squared-return spec that did not match the code that produced these statistics.
 We report two-sided $t$-statistics; as a conservative multiple-testing heuristic inspired by the factor-discovery hurdle of \citet{harvey2016}, we flag whether $|t| > 3.0$, while formal inference rests on the DM/HAC statistic.} confirm that no strategy difference clears the conservative $|t| > 3.0$ heuristic adapted from \citet{harvey2016} (S1 vs.\ S0: $t = 2.42$; S2 vs.\ S0: $t = 0.75$). The limited cross-OOS success rate of 2 in 6 windows suggests that VoV conditioning is sample-dependent. We therefore present S2 as a hypothesis-generating in-sample accounting result rather than a validated trading rule. The two wins are still informative: they occur in windows containing the February 2018 Volmageddon episode and the 2020 COVID crash, consistent with a state-contingent insurance interpretation, but insufficient for contribution-level robustness.

Sensitivity analysis across VoV $z$-score thresholds confirms robustness of the decomposition: at thresholds of 0.5, 1.0, and 1.5, the opportunity cost share remains above 57\% in all cases (64\%, 57\%, and 65\%, respectively), and total premium reduction relative to unconditional VT ranges from 62\% to 76\%. The core finding---that opportunity cost dominates direct cost---is invariant to the specific threshold choice.

%=================================================================
\section{Discussion}
\label{sec:discussion}
%=================================================================

Our decomposition yields three specific implications. First, the conventional view that VT is costly because of high turnover (S1's annualized turnover is 872\%) is misleading. At 5~bps per trade, this turnover generates only 43~bps of annual direct cost---less than a typical ETF expense ratio. The true cost is the 420~bps of foregone equity returns, and this distinction matters for strategy design: reducing trading frequency or using futures addresses at most 9\% of the problem.

Second, the VVIX-conditioning exercise suggests a possible mechanism but does not validate a deployable rule. During low-VoV regimes, the volatility forecast may be accurate but economically unnecessary---remaining fully invested avoids the opportunity cost that drives most of the unconditional VT premium, consistent with \citet{liu2019}, who find that volatility management does not consistently improve performance outside the high-volatility states where timing is most valuable. In the full sample, VoV conditioning behaves like a state-contingent insurance policy and sharply reduces the measured premium. The weak 2-of-6 cross-OOS record, however, means this should be interpreted as an exploratory extension of the decomposition, not as a robust strategy contribution. This connects to the broader ``volatility-of-volatility risk'' literature \citep{bollerslev2009, huang2019, todorov2010, bekaert2014} but applies it to a portfolio construction context rather than asset pricing.

Third, the rebalancing premium quantification reveals a structural benchmark advantage: a 50/50 SPY/GLD portfolio generates a 54~bps full-sample average rebalancing premium annually through mechanical rebalancing, requiring zero forecasting ability (though the premium is negative in two of the four sub-periods). While S2 achieves a higher Sharpe ratio than 50/50 (0.63 vs.\ 0.50), this comparison is not apples-to-apples: S2 is 100\% equity with occasional VT, while 50/50 includes gold diversification that provides crisis alpha without parameter estimation. On a CAGR basis, 50/50 (7.89\%) approaches S1 (7.11\%) with far less complexity, and its long-run average rebalancing premium raises the bar that VT must overcome to deliver net positive alpha---a bar that varies materially across sub-periods.

Several limitations warrant discussion. The VVIX-reliable period begins only in 2012, limiting our sample to 13 years. Cross-OOS analysis now covers all six complete non-overlapping two-year windows, and the 2-of-6 success rate confirms that the specific VoV conditioning rule is not robust enough to stand as an independent trading-rule contribution. The rule ($z > 1.0$ with VIX rising) has not been validated through nested out-of-sample testing; future work should apply walk-forward optimization to confirm whether any related regime-conditioning family is robust. The 5~bps transaction cost assumption is conservative; institutional investors may face lower costs, marginally reducing the direct component. Our analysis focuses on SPY; extension to other asset classes and international markets---particularly those with different volatility dynamics such as emerging markets with amplified VIX sensitivity---is warranted and may reveal different opportunity/direct cost ratios.

%=================================================================
\section{Conclusion}
\label{sec:conclusion}
%=================================================================

We decompose the volatility targeting insurance premium into opportunity cost and direct (transaction) cost components. The key finding---that 91\% of the cost is opportunity cost, not transaction cost---reframes the VT design problem. The path to cheaper VT lies not in lower-friction execution but in understanding when reduced equity exposure is economically worth its opportunity cost. VVIX-conditional targeting reduces cost by 74\% in the full sample, but cross-OOS analysis over all six two-year windows shows only 2 wins against buy-and-hold and no Harvey-significant DM evidence. The structural rebalancing premium embedded in 50/50 SPY/GLD creates a high bar for VT to clear, though VoV-conditional VT achieves a higher full-sample Sharpe ratio (0.63 vs.\ 0.50) via a fundamentally different mechanism (equity timing vs.\ asset-class diversification). We therefore keep the cost decomposition as the robust contribution and treat VoV conditioning as hypothesis-generating: useful for illustrating the opportunity-cost channel, but not yet validated as a standalone trading rule.

%=================================================================
% References
%=================================================================
\bibliographystyle{apalike}
\begin{thebibliography}{99}

\bibitem[Barroso and Santa-Clara(2015)]{barroso2015}
Barroso, P., Santa-Clara, P., 2015. Momentum has its moments. \textit{Journal of Financial Economics} 116(1), 111--120. \url{https://doi.org/10.1016/j.jfineco.2014.11.010}

\bibitem[Bollerslev et al.(2009)]{bollerslev2009}
Bollerslev, T., Tauchen, G., Zhou, H., 2009. Expected stock returns and variance risk premia. \textit{Review of Financial Studies} 22(11), 4463--4492. \url{https://doi.org/10.1093/rfs/hhp008}

\bibitem[Bekaert and Hoerova(2014)]{bekaert2014}
Bekaert, G., Hoerova, M., 2014. The VIX, the variance premium and stock market volatility. \textit{Journal of Econometrics} 183(2), 181--192. \url{https://doi.org/10.1016/j.jeconom.2014.05.008}

\bibitem[Bongaerts et al.(2020)]{bongaerts2020}
Bongaerts, D., Kang, X., van~Dijk, M., 2020. Conditional volatility targeting. \textit{Financial Analysts Journal} 76(4), 54--71. \url{https://doi.org/10.1080/0015198X.2020.1790853}

\bibitem[Todorov(2010)]{todorov2010}
Todorov, V., 2010. Variance risk-premium dynamics: The role of jumps. \textit{Review of Financial Studies} 23(1), 345--383. \url{https://doi.org/10.1093/rfs/hhp035}

\bibitem[CBOE(2012)]{cboe2012}
CBOE, 2012. Double the fun with CBOE's VVIX Index. White Paper, Chicago Board Options Exchange. Current index documentation: \url{https://www.cboe.com/us/indices/dashboard/vvix/}

\bibitem[Booth and Fama(1992)]{booth1992}
Booth, D.G., Fama, E.F., 1992. Diversification returns and asset contributions. \textit{Financial Analysts Journal} 48(3), 26--32. \url{https://doi.org/10.2469/faj.v48.n3.26}

\bibitem[Cederburg et al.(2020)]{cederburg2020}
Cederburg, S., O'Doherty, M.S., Wang, F., Yan, X.S., 2020. On the performance of volatility-managed portfolios. \textit{Journal of Financial Economics} 138(1), 95--117. \url{https://doi.org/10.1016/j.jfineco.2020.04.015}

\bibitem[Fleming et al.(2001)]{fleming2001}
Fleming, J., Kirby, C., Ostdiek, B., 2001. The economic value of volatility timing. \textit{Journal of Finance} 56(1), 329--352. \url{https://doi.org/10.1111/0022-1082.00327}

\bibitem[Harvey et al.(2016)]{harvey2016}
Harvey, C.R., Liu, Y., Zhu, H., 2016. \ldots and the cross-section of expected returns. \textit{Review of Financial Studies} 29(1), 5--68. \url{https://doi.org/10.1093/rfs/hhv059}

\bibitem[Hasbrouck(2009)]{hasbrouck2009}
Hasbrouck, J., 2009. Trading costs and returns for U.S.\ equities: Estimating effective costs from daily data. \textit{Journal of Finance} 64(3), 1445--1477. \url{https://doi.org/10.1111/j.1540-6261.2009.01469.x}

\bibitem[Harvey et al.(2018)]{harvey2018}
Harvey, C.R., Hoyle, E., Korgaonkar, R., Rattray, S., Sargaison, M., Van~Hemert, O., 2018. The impact of volatility targeting. \textit{Journal of Portfolio Management} 45(1), 14--33. \url{https://doi.org/10.3905/jpm.2018.45.1.014}

\bibitem[Hocquard et al.(2013)]{hocquard2013}
Hocquard, A., Ng, S., Papageorgiou, N., 2013. A constant-volatility framework for managing tail risk. \textit{Journal of Portfolio Management} 39(2), 28--40. \url{https://doi.org/10.3905/jpm.2013.39.2.028}

\bibitem[Huang et al.(2019)]{huang2019}
Huang, D., Schlag, C., Shaliastovich, I., Thimme, J., 2019. Volatility-of-volatility risk. \textit{Journal of Financial and Quantitative Analysis} 54(6), 2423--2452. \url{https://doi.org/10.1017/S0022109018001436}

\bibitem[Liu et al.(2019)]{liu2019}
Liu, F., Tang, X., Zhou, G., 2019. Volatility-managed portfolio: Does it really work? \textit{Journal of Portfolio Management} 46(1), 38--51. \url{https://doi.org/10.3905/jpm.2019.1.107}

\bibitem[Moreira and Muir(2017)]{moreira2017}
Moreira, A., Muir, T., 2017. Volatility-managed portfolios. \textit{Journal of Finance} 72(4), 1611--1644. \url{https://doi.org/10.1111/jofi.12513}

\bibitem[Perchet et al.(2016)]{perchet2016}
Perchet, R., de~Carvalho, R.L., Heckel, T., Moulin, P., 2016. Predicting the success of volatility targeting strategies: Application to equities and other asset classes. \textit{Journal of Alternative Investments} 18(3), 21--38. \url{https://doi.org/10.3905/jai.2016.18.3.021}

\end{thebibliography}

\end{document}
